\documentclass[aspectratio=169,10pt]{beamer}

\usepackage[T2A]{fontenc}
\usepackage[utf8]{inputenc}
\usepackage[russian]{babel}
\usepackage{amsmath,amssymb}
\usepackage{graphicx}
\usepackage{booktabs}
\usepackage{hyperref}

\usetheme{Madrid}
\usecolortheme{default}
\setbeamertemplate{navigation symbols}{}
\setbeamertemplate{footline}[frame number]

% Раскомментируйте, чтобы видеть заметки докладчика:
% \usepackage{pgfpages}
% \setbeameroption{show notes on second screen=right}

\title[Слабо информативные априорные распределения]{Weak Informativity and the Information\\ in One Prior Relative to Another}
\subtitle{Доклад по статье M. Evans, G. H. Jang}
\author{Лина Горбунова}
\institute{МФТИ}
\date{\today}

\begin{document}

% ---------- 1. Титул (0:00–0:20) ----------
\begin{frame}
  \titlepage
  \note{Представьтесь, назовите статью и авторов.
  Скажите: работа про то, как формально сравнивать «информативность»
  одного априорного распределения относительно другого — базового.}
\end{frame}

% ---------- 2. План (0:20–0:40) ----------
\begin{frame}{План доклада}
  \begin{enumerate}
    \item Постановка задачи и мотивация
    \item Мера конфликта «априорное — данные» ($P$-значение)
    \item Определение слабой информативности
    \item Примеры: нормальные, $t$, обратно-гамма priors
    \item Применения: биномиальная модель, регрессия, логистическая регрессия
    \item Уточнение через анциллярность и выводы
  \end{enumerate}
  \note{20 секунд. Просто перечислите пункты, не раскрывая.}
\end{frame}

% ---------- 3. Мотивация (0:40–1:40) ----------
\begin{frame}{Мотивация: слабо информативные priors}
  \begin{itemize}
    \item Идея Gelman (2006): между \emph{информативным} и \emph{неинформативным} prior
          вводится компромисс — \alert{слабо информативный} prior.
    \item Пусть $\Pi_1$ — базовый prior, отражающий наши текущие знания о $\theta$.
    \item Хотим выбрать $\Pi_2$, который \emph{консервативен}: вносит меньше информации,
          но \emph{не отбрасывает} $\Pi_1$ полностью.
    \item Причины:
      \begin{itemize}
        \item несобственный $\Pi_2$ создаёт проблемы;
        \item может возникнуть конфликт $\Pi_1$ с данными — хочется его избежать.
      \end{itemize}
    \item \alert{Вопрос:} что значит «$\Pi_2$ менее информативен, чем $\Pi_1$»?
  \end{itemize}
  \note{Главная мысль: интуитивно «шире = менее информативно» не всегда верно
  (см. пример с $t$-prior в разделе 3.2). Нужно строгое определение.}
\end{frame}

% ---------- 4. Обозначения (1:40–2:20) ----------
\begin{frame}{Обозначения и факторизация}
  Модель $\{P_\theta : \theta \in \Theta\}$, $P_\theta(A)=\int_A f_\theta(x)\,\mu(dx)$.
  Prior $\Pi$ даёт \emph{prior predictive}
  \[
    M(A)=\int_\Theta P_\theta(A)\,\Pi(d\theta),\qquad
    m(x)=\int_\Theta f_\theta(x)\,\Pi(d\theta).
  \]
  $T$ — минимальная достаточная статистика. Тогда
  \[
    P_\theta \times \Pi
      = \underbrace{M \times \Pi(\cdot\mid x)}_{\text{выводы}}
      = \underbrace{P(\cdot\mid T)}_{\text{проверка модели}}
        \times \underbrace{M_T}_{\text{конфликт prior--data}}
        \times \Pi(\cdot\mid T).
  \]
  \begin{itemize}
    \item $P(\cdot\mid T)$ — проверка модели;
    \item $M_T$ — проверка \alert{конфликта prior--data}.
  \end{itemize}
  \note{Подчеркните: нас интересует именно компонента $M_T$.
  Она отвечает за «согласованность» prior с данными.}
\end{frame}

% ---------- 5. P-value конфликта (2:20–3:20) ----------
\begin{frame}{$P$-значение конфликта prior--data}
  Наблюдаем $t_0 = T(x_0)$. Сравниваем его с $M_T$:
  \[
    P(t_0) = M_T\bigl(m_T^*(t) \le m_T^*(t_0)\bigr).
    \tag{2--3}
  \]
  \begin{itemize}
    \item $m_T^*$ — инвариантная версия плотности $M_T$
          (учитывает якобиан $J_T(x)$ от $T$).
    \item Идея: если $t_0$ попадает в «хвост» $M_T$, плотность там мала —
          prior и данные конфликтуют.
    \item Асимптотика (Evans, Jang 2011):
          $P(t_0)\to \Pi\bigl(\pi(\theta)\le \pi(\theta_*)\bigr)$ при $n\to\infty$.
    \item При $\Pi_1$ правильном, $P_1(t_0)$ при $t_0\sim M_{1T}$ \alert{равномерно распределено} на $[0,1]$.
  \end{itemize}
  \note{Ключевая идея: $P$-значение — это вероятность «быть в хвосте prior».
  Если она мала — сигнал конфликта.}
\end{frame}

% ---------- 6. Определение слабой информативности (3:20–4:30) ----------
\begin{frame}{Определение слабой информативности}
  Пусть $\Pi_1$ — базовый prior, $t_0 \sim M_{1T}$, $x_\gamma$ — квантиль уровня $\gamma$
  распределения $P_1(t_0)$.
  \[
    M_{1T}\bigl(P_2(t_0) \le x_\gamma\bigr).
    \tag{4}
  \]

  \begin{block}{Определение 1}
    $\Pi_2$ \alert{слабо информативен} относительно $\Pi_1$ на уровне $\gamma$,
    если (4) $\le x_\gamma$.
    \begin{itemize}
      \item Если это верно для всех $\gamma\le\gamma_0$ — \emph{равномерно} на уровне $\gamma_0$.
      \item Если для всех $\gamma$ — \emph{равномерно слабо информативен}.
    \end{itemize}
  \end{block}

  Интерпретация: при использовании $\Pi_2$ (когда верен $\Pi_1$)
  априорная вероятность конфликта prior--data \alert{не выше}, чем у $\Pi_1$.
  \note{Обратите внимание: сравниваем распределение $P_2(t_0)$ при $t_0\sim M_{1T}$.
  Именно это даёт корректное определение.}
\end{frame}

% ---------- 7. Количественная мера (4:30–5:00) ----------
\begin{frame}{Насколько меньше информации?}
  Степень слабой информативности $\Pi_2$ относительно $\Pi_1$:
  \[
    1 - \frac{M_{1T}\bigl(P_2(t_0) \le x_\gamma\bigr)}{x_\gamma}.
    \tag{5}
  \]
  \begin{itemize}
    \item Показывает долю \emph{сокращения} ожидаемых конфликтов prior--data
          при переходе от $\Pi_1$ к $\Pi_2$.
    \item Позволяет выбрать $\Pi_2$, скажем, на $50\%$ менее информативный.
    \item Обычно существует \emph{много} слабо информативных priors —
          нужен дополнительный критерий выбора.
  \end{itemize}
  \note{Это ответ на вопрос «а насколько?». Не просто «да/нет»,
  а количественная характеристика.}
\end{frame}

% ---------- 8. Лемма 1 (5:00–5:40) ----------
\begin{frame}{Упрощённый критерий (Лемма 1)}
  Если $P_i(t)$ непрерывны по $M_{iT}$, то найдётся $r_\gamma$ такое, что
  \[
    M_{1T}\bigl(P_2(t)\le \gamma\bigr) = M_{1T}\bigl(m_{2T}^*(t)\le r_\gamma\bigr).
  \]

  \begin{block}{Эквивалентное условие}
    $\Pi_2$ равномерно слабо информативен относительно $\Pi_1$
    $\iff$
    для всех $t_0$
    \[
      M_{1T}\bigl(m_{2T}^*(t)\le m_{2T}^*(t_0)\bigr)
      \;\le\;
      M_{2T}\bigl(m_{2T}^*(t)\le m_{2T}^*(t_0)\bigr).
    \]
  \end{block}

  \alert{Смысл:} $M_{2T}$ кладёт в «хвосты» плотности $m_{2T}^*$
  не меньше массы, чем $M_{1T}$, т.е. $M_{2T}$ «более размыто».
  \note{Это рабочая форма критерия. Именно ей пользуются в примерах.}
\end{frame}

% ---------- 9. Пример: нормальные priors (5:40–6:30) ----------
\begin{frame}{Пример 1: нормальные priors}
  Модель $x_i\sim N(\mu,1)$, $t=\bar x\sim N(\mu,1/n)$.
  $\Pi_i = N(\mu_0,\sigma_i^2)$, $i=1,2$.

  \[
    M_{iT} = N\!\left(\mu_0,\ \tfrac{1}{n}+\sigma_i^2\right),
  \]
  \[
    M_{1T}\bigl(P_2(t_0)\le \gamma\bigr)
      = 1 - G_1\!\left(
        \frac{1/n+\sigma_2^2}{1/n+\sigma_1^2}\,
        G_1^{-1}(1-\gamma)
      \right).
    \tag{6}
  \]
  \begin{itemize}
    \item (6) $\le\gamma \iff \sigma_2 > \sigma_1$.
    \item При $n\to\infty$ порог переходит в $\sigma_2^2/\sigma_1^2$.
    \item \alert{Обобщение (Теорема 1):} для $\mu\in\mathbb R^k$
          $\Pi_2$ равномерно слабо информативен $\iff$
          $\Sigma_2-\Sigma_1 \succeq 0$.
  \end{itemize}
  \note{Первый важный результат: для нормального prior «менее информативно» = «больше дисперсия».
  Но это лишь частный случай!}
\end{frame}

% ---------- 10. Пример: t vs normal (6:30–7:30) ----------
\begin{frame}{Пример 2: $t$-prior против нормального}
  $\Pi_1 = N(\mu_0,\sigma_1^2)$, $\Pi_2 = t_1(\mu_0,\sigma_2^2,\lambda)$.

  Асимптотический критерий:
  \[
    \frac{\sigma_2^2}{\sigma_1^2}
      \;\ge\;
      K(\lambda)
      = \sup_{\gamma\in[0,1]}
        \frac{G_1^{-1}(1-\gamma)}{H_{1,\lambda}^{-1}(1-\gamma)}.
  \]

  \begin{itemize}
    \item $K(1)=0.6366$: для Cauchy нужно $\sigma_2^2\ge 0.637\,\sigma_1^2$.
    \item $K(3)=0.8488$: $t_1(0,\sigma_2^2,3)$ с \emph{той же дисперсией}, что $N$,
          \alert{не} слабо информативен.
      Нужна дисперсия $\ge 2.55\,\sigma_1^2$.
    \item $K(\lambda)\to 1$ при $\lambda\to\infty$ — совпадает с нормальным случаем.
  \end{itemize}

  \alert{Вывод:} дисперсия — не универсальная мера информативности;
  важна «островершинность» распределения.
  \note{Это ключевой контринтуитивный пример статьи. Сравните с интуицией
  «больше дисперсия = менее информативно» — здесь это не работает.}
\end{frame}

% ---------- 11. Обратно-гамма (7:30–8:00) ----------
\begin{frame}{Пример 3: обратно-гамма priors}
  Модель $N(0,\sigma^2)$, $1/\sigma^2\sim\mathrm{Gamma}_{\mathrm{rate}}(\alpha_i,\beta_i)$.

  Ограничение: совпадение мод $m_{iT}^*$ $\Rightarrow$
  $\beta_2/(\alpha_2+\tfrac12)=\beta_1/(\alpha_1+\tfrac12)$.

  \begin{block}{Теорема 4}
    $\mathrm{Gamma}(\alpha_2,\beta_2)$ асимптотически слабо информативен
    относительно $\mathrm{Gamma}(\alpha_1,\beta_1)$
    $\iff$
    $\alpha_2\le\alpha_1$ и
    $\beta_2 = \beta_1\dfrac{\alpha_2+1/2}{\alpha_1+1/2}$.
  \end{block}

  \alert{Важно:} нельзя неограниченно уменьшать $\beta_2$ — есть нижняя граница.
  \note{Ещё одно предупреждение: наивное «уменьшим параметр» не работает.}
\end{frame}

% ---------- 12. Применения (8:00–9:00) ----------
\begin{frame}{Применения}
  \textbf{Биномиальная модель.} $T\sim\mathrm{Bin}(n,\theta)$, $\theta\sim\mathrm{Beta}(\alpha,\beta)$.
  Для базового $\mathrm{Beta}(6,6)$, $n=20$, $\gamma=0.05$:
  $\mathrm{Beta}(\alpha,\alpha)$ равномерно слабо информативен при $1\le\alpha\le 12.36$.
  При росте $n$ значения $\alpha>6$ отсекаются.

  \vspace{0.3em}
  \textbf{Нормальная регрессия.} $y\sim N(X\beta,\sigma^2 I)$.
  \begin{itemize}
    \item Сначала проверяем prior на $\sigma^2$ по $s^2$;
    \item затем — на $\beta$ по $b\mid s^2$ ($s^2$ анциллярна для $\beta$);
    \item Заменяем $N_k$ на $t_k$: условие $\Sigma_2\succeq\tau_\lambda^2\Sigma_1$.
  \end{itemize}

  \vspace{0.3em}
  \textbf{Логистическая регрессия.} (данные Racine et al.\ 1986; Gelman et al.\ 2008)
  \begin{itemize}
    \item Базовый $\Pi_1=N(0,10^2)\times N(0,2.5^2)$.
    \item \alert{Не любое увеличение} $\sigma_i$ даёт слабую информативность:
          при $\sigma_0\to\infty$ prior predictive $T$ вырождается
          в распределение на двух точках. Есть \emph{конечный} диапазон допустимых $\sigma_i$.
  \end{itemize}
  \note{Логистический пример — самое сильное предостережение статьи.
  Нельзя просто «раздувать» $\sigma$.}
\end{frame}

% ---------- 13. Анcиллярность (9:00–9:30) ----------
\begin{frame}{Уточнение через анциллярность}
  Если $U(T)$ — анциллярная статистика, вариация $U(T)$ не зависит от $\theta$.
  Её нужно удалить из $P$-значения:
  \[
    M_T\bigl(m_T^*(t)\le m_T^*(t_0)\,\big|\,U(T)\bigr).
    \tag{9}
  \]
  Определение (4) заменяется на
  \[
    M_{1T}\bigl(P_2(t_0\mid U(T))\le x_\gamma \,\big|\, U(T)\bigr).
    \tag{10}
  \]
  \begin{itemize}
    \item Если $T$ полная, анциллярности независимы от $T$ (Basu) — conditioning не нужен.
    \item При нескольких максимальных анциллярностях берём \emph{максимум} по ним (10):
          если он мал — уверенно считаем $\Pi_2$ слабо информативным.
  \end{itemize}
  \note{Техническое, но важное уточнение: для корректного $P$-значения
  нужно conditioning на анциллярность.}
\end{frame}

% ---------- 14. Выводы (9:30–10:00) ----------
\begin{frame}{Выводы}
  \begin{itemize}
    \item Предложен \alert{формальный} критерий информативности prior
          относительно базового — через вероятность конфликта prior--data.
    \item Критерий даёт \alert{количественную} меру (5):
          на сколько процентов меньше конфликтов ожидать.
    \item В типичных семействах (нормальные, $t$, обратно-гамма)
          результаты интуитивно разумны,
          но \alert{дисперсия — не универсальная мера}.
    \item Логистическая регрессия: рост $\sigma$ может \emph{ухудшить}
          слабую информативность — есть конечный допустимый диапазон.
    \item Открытые вопросы: связь с неинформативными/несобственными priors,
          выбор среди множества слабо информативных priors.
  \end{itemize}

  \vspace{0.5em}
  \centering
  \alert{Спасибо за внимание!}
  \note{Завершите: подчеркните, что статья даёт рабочее определение
  и предупреждает о ловушках. Готовы к вопросам.}
\end{frame}

\end{document}